
Рассмотрим подробнее структуру характеристического уравнения. Представим его в виде слагаемого из произведения диагональных элементов матрицы и многочлена $Q$, где $Q$ состоит из всех прочих слагаемых из явной формулы определителя (в каждое из них входит не более чем $n-2$ скобки вида $(a_{ij}-\lambda)$). Значит $\deg Q(\lambda) \leqslant n-2$.
\[
\det (A-\lambda E)=\det
        \begin{pmatrix}
        a_{11}-\lambda & a_{12} & \cdots & a_{1n} \\
        a_{21} & a_{22}-\lambda & \cdots & 0 \\
        \vdots & \vdots & \ddots & \vdots \\
        a_{n1} & 0 & \cdots & a_{nn}-\lambda
        \end{pmatrix}=(a_{11} - \lambda)\dots(a_{nn}-\lambda)+Q(\lambda)
\]
Тогда если рассмотреть характеристическое уравнение как характеристический многочлен $P(\lambda)$, то его можно представить в виде:
\[
P(\lambda)=(-1)^n\lambda^n+(-1)^{n-1}\lambda^{n-1}(a_{11}+\dots+a_{nn})+\dots+\det A
\]
Коэффициент $a_{11}+\dots+a_{nn}$ называется следом матрицы $A$ и  обозначается $tr\ A$.
\begin{theorem}
    Если $A$ и $A'$ матрицы преобразования $\phi$ в различных базисах, то характеристические многочлены матриц $A$ и $A'$ совпадают.
\end{theorem}
\begin{proof}
    Пусть $S$ матрица перехода между базисами. Тогда, согласно ранее полученной \hyperref[A']{формуле}: $A'-\lambda E \to S^{-1}(A-\lambda E)S$. Следовательно, $\det (A'-\lambda E)=\det (S^{-1}(A-\lambda E)S)=\det S^{-1} \det (A-\lambda E) \det S=\det (A-\lambda E)$. Откуда получаем требуемое $\det (A'-\lambda E)=\det (A-\lambda E)$.
\end{proof}
Из этой теоремы мы понимаем, что характеристический многочлен не зависит от базиса. Значит, далее корректно будет говорить о том, что это \textit{характеристический многочлен преобразования} (а не матрицы преобразования). Так же можно зафиксировать следующее следствие: $\det A=\det A'$ и $tr \ A=tr\ A'$.
\begin{proposition}
    Собственные векторы преобразования $\phi$, отвечающие различным собственным значениям, образуют ЛНЗ набор.
\end{proposition}
\begin{proof}
    Пусть $\{\lambda_i|i\in \overline{1,m}\}$ -- собственные значения, а $\{h_i|i\in \overline{1,m}\}$ -- соответствующие собственные векторы. Рассмотрим вспомогательное преобразование $B_i=A - \lambda E$. Тогда \[B_i(h_j)=Ah_j-\lambda_i h_j=\lambda_j h_j - \lambda_i h_j=(\lambda_j - \lambda_i)h_j = \begin{cases} 
   0, \ \ i=j \\
   \neq0, \ \ i \neq j 
\end{cases}\]
Предположим, что $h_1, \dots, h_m$ - ЛЗ. Тогда не умаляя общности $h_1=\alpha_2 h_2+\dots+\alpha_m h_m$. Применим к этому равенству композицию $B_mB_{m-1}\dots B_2$. Тогда $B_m\dots B_2(h_1)=(\lambda_1-\lambda_m)\dots (\lambda_1-\lambda_2)h_1\neq0$. Но $\forall i\in \overline{1,m}\hookrightarrow B_m\dots B_2(h_i)=0$. Противоречие.
\end{proof}
\textit{Следствие}: собственные подпространства образуют прямую сумму.

\textit{Напоминание}: пусть $P(\lambda)=(\lambda-\lambda_0)^SQ(\lambda)$, где $S$ максимально возможное, тогда $\lambda_0$ -- корень кратности $S$.
\begin{theorem}
    Пусть $\lambda_0$ корень кратности $S$ характеристического многочлена преобразования $A:L\to L$. Тогда размерность $t$ собственного подпространства $L'$, соответствующего $\lambda_0$, удовлетворяет неравенству $1\leqslant t\leqslant S$.
\end{theorem}
\begin{proof}
    Выберем в $L'$ базис $e_1, \dots, e_t$ и дополним его до базиса в $L$ векторами $e_{t+1}, \dots, e_n$. Тогда в этом базисе
\[
A=\left(
\begin{array}{c|c}
\begin{matrix}
\lambda_0 & \cdots & 0 \\
\vdots & \ddots& \vdots\\
0 & \cdots& \lambda_0
\end{matrix}
& 
\begin{matrix}
B
\end{matrix} \\
\hline
\begin{matrix}
0
\end{matrix}
& 
\begin{matrix}
C
\end{matrix}
\end{array}
\right), \ \ \ A(e_i)=\lambda_0 e_i \ \ \forall i \in \overline{1,t}
\]
А значит $\det (A-\lambda E)=(\lambda-\lambda_0)^t\tilde{Q}(\lambda)=P(\lambda)=(\lambda-\lambda_0)^SQ(\lambda)$. Следовательно $t\leqslant S$. Вторая часть неравенства следует из определения корня (любой корень имеет кратность хотя бы 1).
\end{proof}
Иногда $S$ называют \textit{алгебраической кратностью}, а $t$ \textit{геометрической кратностью}.
\textit{Замечание}: $S$ не всегда совпадает с $t$. Пример, когда $S\neq t$:
преобразование с матрицей $A=\left(\begin{matrix}
    1&1\\0&1
\end{matrix}\right)$. Тогда $\lambda=1$ -- корень кратности $S=2$, но геометрическая кратность $t=1$: \[
\det (A-\lambda E)=(1-\lambda)^2=0, \ \ \ A-\lambda E=\left(\begin{matrix}
    0&1\\0&0
\end{matrix}\right) \to t=1
\]

\textit{Напоминание}: пусть $P(\lambda)$ -- многочлен с вещественными коэффициентами; если $\lambda_i \in \mathbb{C}$ его корень, то $\overline\lambda$ тоже корень ($\overline\lambda$ - сопряженное с $\lambda$).
\begin{theorem}
    Паре сопряженных корней $\lambda, \overline{\lambda}$ характеристического многочлена преобразования $A$ будет соответствовать ненулевое инвариантное относительно $A$ подпространство $L'$ такое, что оно не содержит собственных векторов и через любой его ненулевой вектор проходит 2-мерное инвариантное подпространство.
\end{theorem}
\begin{proof}
    Пусть $t^2+pt+q=0$ квадратное уравнение с корнями $\lambda, \overline{\lambda}$. Рассмотрим вспомогательное преобразование $B=A^2+pA+qE$. Возьмем $L'=\ker B$. Тогда $L'$ инвариантно относительно $A$.
    \begin{itemize}
        \item Преобразуем $B=A^2+pA+qE=(A-\lambda E)(A-\overline{\lambda}E)$. Так как $\lambda$ корень характеристического многочлена, то $\det (A-\lambda E)=0$, тогда $\det B=0$. Следовательно, $L'$ непустое.
        \item Предположим, что $h\in L'$ -- собственный вектор, тогда $Ah=\mu h$ и $Bh=0$. Но $Bh=A(h)^2+pA(h)+qE=(\mu^2+p\mu+q)h$, а $\mu^2+p\mu+q \neq 0$, противоречие.
    \end{itemize}
    Теперь рассмотрим линейную оболочку векторов $x\in L'$ и $A(x)\in L'$ (они не коллинеарны, так как иначе $x$ -- собственный). Эта оболочка образует инвариантное подпространство $L''\subset L'$ размерности 2:
    \[A(\alpha x+\beta A(x))=\alpha A(x)+\beta A^2(x)\in L''\]
    Следовательно $L'$ удовлетворяет условию теоремы.
\end{proof}

\section{Диагонализируемость матриц}
\begin{theorem}
    Матрица преобразования $A$ в базисе $e^{(n)}$ имеет диагональный вид тогда и только тогда, когда $e_1, \dots, e_n$ -- собственные векторы.
\end{theorem}
\begin{proof} \textit{(не лекторское)}
    \begin{itemize}
        \item[$\Rightarrow$] Если матрица $A$ диагональная, то $Ae_i=a_{ii}e_i$, откуда $e_i$ -- собственный.
        \item [$\Leftarrow$] Если векторы $e_1, \dots, e_n$ собственные, то $\forall i\in \overline{1,n} \hookrightarrow A(e_i)=\lambda_i e_i$, откуда матрица $A$ -- диагональная.
    \end{itemize}
\end{proof}
